\chapter{Bộ đếm dùng chung cho các cửa sổ chồng lấn}
\chapterlead{Khi một ràng buộc trượt trên chuỗi, hai cửa sổ kề nhau chia sẻ gần
toàn bộ biến. Mã hóa từng cửa sổ độc lập sẽ lặp trạng thái; bộ đếm dùng chung
(\emph{shared counter}) biến phần giao đó thành một mô-đun có thể tái sử dụng.}

\index{SEQUENCE constraint}\index{shared counter}
\index{Sequential Counter Ladder}\index{Ladder AMK}
\index{sliding window}

\flowdiagram{Chuỗi Boolean}{Phân hoạch khối}{Bộ đếm cục bộ}
  {Bộ nối}{Họ cửa sổ}

\section{SEQUENCE trong văn liệu về ràng buộc}

Ràng buộc cửa sổ thuộc một họ lâu đời hơn các bộ đếm dùng chung.
\textsc{Sequence}
yêu cầu mọi đoạn liên tiếp độ dài cố định có số biến nhận một giá trị trong
khoảng \([\ell,u]\). Nó xuất hiện trong xếp thứ tự xe, lập lịch ca, xác định
cỡ lô và các mô hình chuỗi. Vấn đề cốt lõi không chỉ là số literal đúng trong
một cửa sổ, mà còn là sự phụ thuộc giữa toàn bộ các cửa sổ chồng lấn.

Văn liệu đã phát triển ít nhất bốn cách nhìn:
\begin{description}[style=nextline]
  \item[Phân rã theo cửa sổ]
  Tạo một ràng buộc đếm cho mỗi cửa sổ. Cách này có tính mô-đun và dễ dùng,
  nhưng nếu các trạng thái phụ không được chia sẻ thì phần giao bị biểu diễn
  nhiều lần.

  \item[Tổng tiền tố]
  Đặt \(S_i=\sum_{j\le i}x_j\), khi đó mỗi cửa sổ là
  \(\ell\le S_{i+w}-S_i\le u\). Mọi cửa sổ dùng chung một dãy trạng thái tiền tố.

  \item[Ôtômát/ràng buộc chính quy]
  Trạng thái ghi đủ lịch sử cần thiết để quyết định ký hiệu tiếp theo. Một
  ràng buộc \textsc{Regular} có thể biểu diễn việc chuỗi có thuộc một ngôn ngữ
  hữu hạn trạng thái hay không \parencite{pesant2004regular}.

  \item[Luồng hoặc đồ thị]
  Biểu diễn các hiệu tiền tố và cận cửa sổ bằng mạng luồng, từ đó xây dựng
  bộ lan truyền toàn cục.
\end{description}

So sánh có hệ thống các phép mã hóa \textsc{Sequence} cho thấy cách phân rã,
ôtômát và phương pháp dựa trên tổng khác nhau rõ về kích thước và mức nhất quán
\parencite{brand2007sequence}. Các bộ lan truyền dựa trên luồng tiếp tục khai
thác cấu trúc toàn cục để đạt khả năng lọc mạnh \parencite{maher2008sequence};
những khung tổng quát hơn còn xem ràng buộc trên chuỗi như tổ hợp của một
\emph{chữ ký} và một \emph{đặc trưng} cần đếm
\parencite{monette2014sequences}.

\begin{center}
\captionof{table}{Các cách tái sử dụng trạng thái trong ràng buộc cửa sổ.}
\label{tab:window-reuse}
\begin{tabularx}{\textwidth}{@{}>{\bfseries\sffamily}p{.19\textwidth}YYY@{}}
\toprule
Họ & Trạng thái dùng chung & Điểm mạnh & Chi phí/rủi ro\\
\midrule
Mỗi cửa sổ & Không hoặc ít & Đơn giản, dùng bộ mã hóa sẵn & Lặp trạng thái phụ\\
Tổng tiền tố & Tổng từ đầu chuỗi & Mọi cửa sổ là hiệu hai tiền tố & Miền tổng toàn cục\\
Ôtômát & Hậu tố/lịch sử & Xử lý mẫu phức tạp & Bùng nổ trạng thái theo \(w\)\\
Luồng & Nút và cung & Lọc toàn cục & Bộ lan truyền/phép mã hóa phức tạp\\
Bộ đếm theo khối & Tiền tố/hậu tố trong khối & Cục bộ và tái sử dụng có kiểm soát &
Cần chọn phân hoạch\\
\bottomrule
\end{tabularx}
\end{center}

Bộ đếm dùng chung theo khối của chương này nằm giữa hai cực. Nó không tạo một
trạng thái toàn cục cho mọi tiền tố, nhưng cũng không xây bộ đếm độc lập cho
từng cửa sổ. Phân hoạch quyết định phạm vi tái sử dụng, còn bộ nối giữ ngữ
nghĩa của từng cửa sổ.

\section{Họ ràng buộc cửa sổ}

Cho chuỗi
\[
  \Omega=\langle x_1,\ldots,x_n\rangle.
\]
Cửa sổ độ rộng \(w\) bắt đầu tại \(i\):
\[
  W_i(w)=\langle x_i,\ldots,x_{i+w-1}\rangle.
\]
AMK bậc thang là họ
\[
  \sum_{x\in W_i(w)}x\le k
  \qquad i=1,\ldots,n-w+1.
\]
Các ràng buộc chuỗi tổng quát cho phép cả cận dưới và cận trên:
\[
  \ell\le
  \sum_{x\in W_i(w)}x
  \le u.
\]

Nếu mỗi cửa sổ dùng một bộ đếm tuần tự riêng, kích thước xấp xỉ
\[
  O((n-w+1)wk).
\]
Trong khi đó, hai cửa sổ kề nhau chỉ khác một biến đi ra và một biến đi vào.
Mục tiêu của phép mã hóa dùng chung là đưa chi phí về gần
\[
  O(nk)
\]
hoặc \(O(nk+n_c)\), với \(n_c\) là số bộ nối.

\begin{figure}[tb]
\centering
\begin{tikzpicture}[satfig,x=12mm,y=7mm]
  \foreach \i in {1,...,8}{
    \node[satbox,minimum width=10mm,minimum height=8mm] (x\i) at (\i,0)
      {\(x_{\i}\)};
  }
  \draw[Blue,very thick] (0.55,-.8)--(4.45,-.8)
    node[midway,below,satlabel,text=Blue]{\(W_1=\{x_1,\ldots,x_4\}\)};
  \draw[Teal,very thick] (1.55,-1.65)--(5.45,-1.65)
    node[midway,below,satlabel,text=Teal]{\(W_2=\{x_2,\ldots,x_5\}\)};
  \draw[Gold,very thick] (2.55,-2.5)--(6.45,-2.5)
    node[midway,below,satlabel,text=Gold]{\(W_3=\{x_3,\ldots,x_6\}\)};
  \draw[satdim] (2.05,.72)--(4.95,.72)
    node[midway,above,satlabel]{phần trạng thái có thể dùng lại};
\end{tikzpicture}
\caption{Ba cửa sổ độ rộng \(4\) chỉ dịch một vị trí: cửa sổ kề nhau giữ lại
ba trong bốn biến. Xây lại một bộ đếm cho từng hàng sẽ lặp phần giao này.}
\label{fig:overlapping-windows}
\end{figure}

\section{Phân hoạch thành khối}

Chia chuỗi thành các khối liên tiếp
\[
  B_1,B_2,\ldots,B_m,
  \qquad
  \Omega=B_1\cdot B_2\cdots B_m.
\]
Trong mỗi khối, xây hai họ biến ngưỡng:
\[
  F^b_{i,q}=1
  \iff
  \sum_{j=1}^{i}B_b[j]\ge q,
\]
\[
  R^b_{i,q}=1
  \iff
  \sum_{j=i}^{\card{B_b}}B_b[j]\ge q.
\]
\(F\) đếm tiền tố từ trái sang phải; \(R\) đếm hậu tố từ phải sang trái. Một
cửa sổ cắt biên giữa \(B_b,B_{b+1}\) là hợp của một hậu tố bên trái và một
tiền tố bên phải. Số literal đúng trong cửa sổ được đọc từ \(R^b\) và
\(F^{b+1}\).

\begin{figure}[tb]
\centering
\begin{tikzpicture}[satfig,x=11mm,y=9mm]
  \foreach \i in {1,...,8}{
    \node[satbox,minimum width=9mm] (x\i) at (\i,0) {\(x_{\i}\)};
  }
  \node[satlabel,text=Blue] at (2.5,.75) {\(B_1\)};
  \node[satlabel,text=Teal] at (6.5,.75) {\(B_2\)};
  \draw[Blue,thick,rounded corners=2pt] (.55,-.48) rectangle (4.45,.48);
  \draw[Teal,thick,rounded corners=2pt] (4.55,-.48) rectangle (8.45,.48);
  \draw[Gold,ultra thick] (2.55,-.75)--(6.45,-.75)
    node[midway,below,satlabel,text=Gold]{cửa sổ \(W_3=x_3,\ldots,x_6\)};
  \node[satstate] (r) at (3,-2.25) {\(R^1\)};
  \node[satstate] (f) at (6,-2.25) {\(F^2\)};
  \node[satbox,fill=PaleGold,draw=Gold] (c) at (4.5,-3.55)
    {bộ nối\\\(\neg R^1_{3,q}\lor\neg F^2_{2,k-q+1}\)};
  \draw[satedge] (x3.south)--(r);
  \draw[satedge] (x4.south)--(r);
  \draw[satedge] (x5.south)--(f);
  \draw[satedge] (x6.south)--(f);
  \draw[satedge] (r)--(c);
  \draw[satedge] (f)--(c);
\end{tikzpicture}
\caption{Một cửa sổ cắt biên được phân rã thành hậu tố của khối trái và tiền tố
của khối phải; bộ nối chỉ đọc các ngưỡng cần thiết.}
\label{fig:block-connector}
\end{figure}

\begin{example}
Giả sử cửa sổ gồm hậu tố \(A\) của khối trái và tiền tố \(B\) của khối phải.
Đặt \(T(S,q)\) là literal ``\(\sum_{x\in S}x\ge q\)'', với hai quy ước biên
\(T(S,q)=\top\) khi \(q\le 0\) và \(T(S,q)=\bot\) khi \(q>\card{S}\).
Điều kiện at-most-\(k\) tương đương với việc không tồn tại
\(q\in\set{0,\ldots,k+1}\) sao cho
\[
  \sum_{x\in A}x\ge q
  \quad\text{và}\quad
  \sum_{x\in B}x\ge k+1-q.
\]
Bộ nối vì thế chứa
\[
  \neg T(A,q)
  \lor
  \neg T(B,k+1-q),
  \qquad q=0,\ldots,k+1.
\]
Hai giá trị biên \(q=0\) và \(q=k+1\) không phải chi tiết hình thức: chúng
phát hiện trường hợp toàn bộ \(k+1\) giá trị đúng nằm ở chỉ một phía của biên.
\end{example}

\begin{workedexample}{Đọc một cửa sổ qua hai thanh ghi}
Lấy \(n=8,w=4,k=2\), chia
\(B_1=(x_1,\ldots,x_4)\), \(B_2=(x_5,\ldots,x_8)\), và xét
\[
  x=(1,0,1,0\mid 1,0,0,1).
\]
Cửa sổ \(W_3=(x_3,x_4\mid x_5,x_6)=(1,0\mid1,0)\) có tổng \(2\).
Hậu tố trái đạt ngưỡng \(1\) nhưng không đạt \(2\):
\(R^1_{3,1}=1,R^1_{3,2}=0\). Tiền tố phải cũng có
\(F^2_{2,1}=1,F^2_{2,2}=0\). Với \(k=2\), bốn bộ nối ứng với
\(q=0,1,2,3\) cấm các cặp ngưỡng
\((\top,F_3),(R_1,F_2),(R_2,F_1),(R_3,\top)\).
Tất cả đều thỏa: hai mệnh đề biên xác nhận mỗi phía riêng lẻ không đạt \(3\),
còn ở hai mệnh đề giữa luôn có một ngưỡng bằng \(0\).

Nếu đổi \(x_6\) thành \(1\), tiền tố phải đạt ngưỡng \(2\). Khi đó
\(\neg R^1_{3,1}\lor\neg F^2_{2,2}\) sai, phát hiện ngay tổng cửa sổ bằng
\(3\). Ví dụ cho thấy bộ nối không cộng lại toàn cửa sổ; nó chỉ ghép hai
chứng cứ ngưỡng.
\end{workedexample}

\section{Bộ nối như một giao diện}

Bộ đếm cục bộ không biết cửa sổ nào sẽ đọc nó. Nó chỉ cung cấp đặc tả của các
literal ngưỡng. Bộ nối biết cách một ràng buộc cắt biên và chọn các ngưỡng cần
ghép.

Với at-least-\(k\):
\[
  \sum_{x\in A\cup B}x\ge k
\iff
  \bigwedge_{q=1}^{k}
  \left[
    \left(\sum_{x\in A}x\ge q\right)
    \lor
    \left(\sum_{x\in B}x\ge k-q+1\right)
  \right].
\]
Nếu thanh ghi có đặc tả tương đương, mỗi ngoặc vuông trở thành
\[
  R^b_{a,q}\lor F^{b+1}_{c,k-q+1}.
\]

\begin{designrule}{Bên cung cấp và bên sử dụng}
Bộ đếm là bên cung cấp literal ngưỡng; bộ nối là bên sử dụng. Đặc tả hai
chiều cần thiết khi bộ nối dùng cả thanh ghi dương và phủ định để diễn giải
``đạt'' và ``chưa đạt'' ngưỡng.
\end{designrule}

\section{Staircase AMO và SCL}

Với \(k=1\), mỗi cửa sổ thỏa AMO. Thang bộ đếm dùng chung
(\emph{Shared Counter Ladder}, SCL) xây các trạng thái AMO/AMZ trong khối và
nối chúng qua các biên. So với việc tạo xung đột từng cặp trong từng cửa sổ,
SCL không lặp các cặp nằm trong phần giao.

Giả sử \(m=\lceil n/w\rceil\) khối có kích thước gần \(w\). Chi phí bộ đếm cục
bộ tuyến tính theo tổng kích thước khối, còn các bộ nối tuyến tính theo số biên và
ngưỡng. Vì vậy kích thước toàn họ có thể giảm từ bậc hai về tuyến tính theo
\(n\) ở các AMO bậc thang điển hình.

Mã hóa Duplex dùng BDD tiến/lùi trong mỗi cửa sổ và liên kết các trạng thái
\parencite{fazekas2020duplex}. SCL dùng trạng thái tuần tự, nhờ đó giảm số
biến phụ và mệnh đề trong phân tích của
\textcite{truong2025scamo}.

Trong hệ phân loại trên, SCL phục vụ các tập xung đột có dạng khoảng dịch
chuyển và phần giao lớn giữa các khoảng kề nhau. Mã hóa từng cặp và bộ đếm tuần
tự độc lập phù hợp hơn với các tập AMO rời rạc hoặc thiếu một thứ tự tuyến tính
ổn định. Độ chồng lấn vì thế là đặc trưng đầu tiên để chọn giữa trạng thái dùng
chung và bộ đếm cục bộ.

\begin{resultbox}{SCAMO và Anti-Bandwidth}
Trên họ SCAMO với \(n=1000\) và độ rộng cửa sổ thay đổi,
\textcite{truong2025scamo} báo SCL dùng ít biến và mệnh đề hơn Duplex cùng các
phép mã hóa từng cửa sổ. Khi đưa vào bài toán Anti-Bandwidth, SCL đạt kết quả
cạnh tranh trên 24 đồ thị kiểm chuẩn, đặc biệt ở các bài toán lớn nơi chi phí
lặp cửa sổ trở nên rõ rệt.
\end{resultbox}

\section{Mở rộng sang Ladder AMK}

Khi \(k>1\), mỗi khối cần nhiều mức ngưỡng. Ta dùng bộ đếm tuần tự:
\[
  F^b_{i,q}
  \leftrightarrow
  F^b_{i-1,q}\lor
  (B_b[i]\land F^b_{i-1,q-1}),
\]
và hệ thức truy hồi đối xứng cho \(R^b\).

Bộ nối cho cận trên cấm mọi cặp ngưỡng có tổng \(k+1\):
\[
  \neg T(A,q)
  \lor
  \neg T(B,k+1-q),
  \qquad q=0,\ldots,k+1.
\]
Vì thế các thanh ghi thực phải được xây đến ngưỡng \(k+1\), trừ những mức có
thể giản lược thành \(\top\) hoặc \(\bot\) theo kích thước phần cửa sổ.
Nếu một cửa sổ đi qua nhiều hơn một biên, có hai lựa chọn:
\begin{enumerate}
  \item tăng độ rộng khối để mỗi cửa sổ cắt nhiều nhất một biên;
  \item xây bộ nối nhiều ngôi hoặc một bộ đếm trung gian trên tổng của các khối.
\end{enumerate}
Thiết kế thứ nhất đơn giản hơn; thiết kế thứ hai linh hoạt với nhiều width.

Ladder AMK trong \textcite{truong2026ladderamk} phát triển bộ đếm tuần tự dùng
chung cho một tập ràng buộc đếm dạng bậc thang. Giá trị cốt lõi không nằm ở
một mệnh đề riêng lẻ mà ở việc bộ đếm được tạo cho cả họ ràng buộc.

\subsection{Quan hệ với tổng tiền tố và ràng buộc chuỗi tổng quát}

Ladder AMK có thể được mô tả bằng ma trận liên thuộc: mỗi hàng là một ràng
buộc, mỗi cột là một đầu vào, và các giá trị \(1\) trong mỗi hàng tạo thành
một khoảng; các biên khoảng thay đổi đơn điệu theo hàng. Mã hóa tổng tiền tố
dùng một trạng thái toàn cục cho mọi cột. Ladder AMK chỉ vật hóa những ngưỡng
tiền tố/hậu tố mà các hàng thật sự đọc. Do đó tiêu chí lựa chọn không chỉ là
\(n,k\), mà còn là độ phức tạp của hai đường biên bậc thang.

Với nhiều độ rộng và nhiều cận, một ôtômát có thể nắm toàn bộ lịch sử nhưng số
trạng thái tăng theo tổ hợp các bộ đếm đang hoạt động. Số ngưỡng dùng chung tăng
theo số biên và mức đếm được yêu cầu. Hai hướng vì thế tạo một đánh đổi:
ôtômát hợp nhất các ràng buộc trong trạng thái, còn Ladder AMK giữ các ràng
buộc tách biệt nhưng dùng chung đại lượng trung gian. Đây là câu hỏi thực
nghiệm và lý thuyết đáng quan tâm hơn việc chỉ so số mệnh đề trên một cấu hình.

\section{Đúng đắn}

\begin{theorem}[Tính đúng của phép mã hóa cửa sổ dùng chung]
\label{thm:shared-window-correct}
Giả sử:
\begin{enumerate}
  \item các thanh ghi tiến/lùi đúng theo đặc tả ngưỡng;
  \item mọi cửa sổ nằm trong một khối được ràng buộc bởi bộ đếm cục bộ;
  \item mỗi cửa sổ hoặc nằm trong một khối, hoặc cắt đúng một biên khối;
  \item với cửa sổ \(A\cup B\) cắt biên, phép mã hóa chứa đủ các mệnh đề
  \[
    \neg T(A,q)\lor\neg T(B,k+1-q),
    \qquad q=0,\ldots,k+1.
  \]
\end{enumerate}
Khi đó \CNF tương đương theo phép chiếu với họ staircase AMK.
\end{theorem}

\begin{proof}
Để chứng minh tính đúng, xét một cửa sổ. Nếu nằm trong khối, bộ đếm cục bộ loại mọi tổng
vượt \(k\). Nếu cắt biên, đặt \(a=\sum_{x\in A}x\),
\(b=\sum_{x\in B}x\). Khi \(a+b\ge k+1\), chọn
\(q=\min\{a,k+1\}\). Ta có \(T(A,q)=1\); nếu \(q<k+1\) thì
\(b\ge k+1-a=k+1-q\), còn nếu \(q=k+1\) thì ngưỡng bên phải bằng \(0\)
và đúng theo quy ước. Bộ nối tương ứng bị vi phạm.

Để chứng minh tính đầy đủ, lấy một chuỗi thỏa mọi cửa sổ và gán các thanh ghi
bằng đúng tổng tiền tố/hậu tố. Các mệnh đề của bộ đếm đều đúng. Nếu một bộ nối sai thì tồn
tại \(q\) sao cho \(a\ge q\) và \(b\ge k+1-q\), suy ra \(a+b\ge k+1\),
mâu thuẫn với giả thiết cửa sổ nhiều nhất \(k\). Do đó phép gán trên biến gốc
luôn mở rộng được thành mô hình của \CNF; kết hợp với tính đúng cho tương đương
theo phép chiếu.
\end{proof}

\subsection{Cửa sổ cắt nhiều biên}

Định lý~\ref{thm:shared-window-correct} bao phủ trực tiếp cửa sổ trong một khối
và cửa sổ cắt một biên. Với cửa sổ đi qua ba khối trở lên, tổng được ghép từ
nhiều hơn hai thành phần.
Có thể ghép tuần tự các khối bằng một bộ đếm trung gian, hoặc dùng một cây ghép
có đầu ra bị cắt ở \(k+1\). Với các tổng bộ phận
\(S_1,\ldots,S_r\), bất biến cần giữ là
\[
  U_{j,q}\iff \min\!\left(k+1,\sum_{h=1}^{j}S_h\right)\ge q .
\]
Sau lần ghép cuối, mệnh đề \(\neg U_{r,k+1}\) đóng AMK. Chi phí tăng theo số
khối mà từng cửa sổ chạm tới; một cây ghép riêng cho mỗi cửa sổ làm giảm tỷ
lệ trạng thái dùng chung. Điều kiện ``mỗi cửa sổ cắt nhiều nhất một biên'' do
đó là tiêu chí phân hoạch, còn bộ đếm ghép là phương án cho phân hoạch tổng
quát.

\section{Chọn độ rộng khối}

Khối nhỏ tạo nhiều biên và bộ nối nhưng bộ đếm cục bộ ngắn. Khối lớn giảm số
biên nhưng giữ nhiều thanh ghi không được mọi cửa sổ dùng. Vì vậy độ rộng là
một siêu tham số cấu trúc.

Ba cách chọn:
\begin{itemize}
  \item bằng độ rộng cửa sổ khi mọi ràng buộc có cùng \(w\);
  \item theo ước chung lớn nhất hoặc cụm của nhiều độ rộng;
  \item phân hoạch thích nghi, đặt biên tại nơi có ít cửa sổ cắt qua.
\end{itemize}

Phân hoạch thích nghi có thể xem là bài toán tối ưu trên đồ thị khoảng: mỗi cửa
sổ là một khoảng, chi phí của một biên bằng số ràng buộc cắt biên nhân với số
ngưỡng cần nối.

\section{Triển khai}

Một giao diện lập trình (API) cho mô-đun có thể là:
\[
  \texttt{SharedCounter(sequence, partition, maxThreshold)}
\]
và trả về:
\begin{itemize}
  \item các literal ngưỡng tiền tố/hậu tố;
  \item hàm \texttt{prefix(block,pos,q)};
  \item hàm \texttt{suffix(block,pos,q)};
  \item ánh xạ thanh ghi để gỡ lỗi và kiểm tra nội bộ.
\end{itemize}
Sau đó \texttt{addAtMostWindow(i,w,k)} chỉ sinh các bộ nối. Nhiều mô-đun sử
dụng có thể đọc cùng một bộ đếm: AMK, ALK, EXK hoặc hàm mục tiêu.

\begin{algorithm}[H]
\caption{Sinh AMK cửa sổ dùng chung khi mỗi cửa sổ cắt nhiều nhất một biên}
\label{alg:shared-window}
\begin{minipage}{.94\textwidth}\small
\textbf{Đầu vào:} chuỗi \(x_1,\ldots,x_n\); một phân hoạch khối; tập cửa sổ
\(\mathcal W\); cận \(k(W)\) của từng cửa sổ.\\
\textbf{Đầu ra:} một \CNF tương đương theo phép chiếu.
\begin{enumerate}
  \item Trong mỗi khối, sinh các thanh ghi ngưỡng tiến/lùi đến mức lớn
  nhất \(\max_{W\in\mathcal W}(k(W)+1)\) thật sự được đọc.
  \item Với cửa sổ nằm trọn trong một khối, đọc bộ đếm cục bộ và cấm
  ngưỡng \(k(W)+1\).
  \item Với cửa sổ \(W=A\cup B\) cắt một biên, lần lượt
  \(q=0,\ldots,k(W)+1\), sinh
  \(\neg T(A,q)\lor\neg T(B,k(W)+1-q)\); giản lược các literal hằng.
  \item Nếu phát hiện một cửa sổ cắt từ hai biên trở lên, dừng mô-đun này
  và chuyển sang bộ đếm ghép; không áp dụng bộ nối hai ngôi.
  \item Trả \CNF cùng ánh xạ từ mỗi thanh ghi về
        \((\text{khối},\text{vị trí},q)\).
\end{enumerate}
\end{minipage}
\end{algorithm}

\subsection{Miền hiệu quả của bộ đếm dùng chung}

Bộ đếm dùng chung đạt lợi thế rõ nhất khi các cửa sổ chồng lấn mạnh, cận \(k\)
nhỏ và mỗi cửa sổ chỉ cắt ít khối. Với các cửa sổ gần rời nhau hoặc bài toán
nhỏ, mã hóa từng cặp và bộ đếm độc lập thường có hằng số thấp hơn. Một bộ chọn
có thể quyết định trước khi sinh \CNF từ ba đại lượng: tổng số ngưỡng cục bộ,
số bộ nối sau giản lược và số bước ghép của các cửa sổ nhiều biên.

\section{Bài học thiết kế}

\begin{summarybox}
\begin{enumerate}
  \item Cửa sổ chồng lấn là tín hiệu nên dùng trạng thái chung.
  \item Bộ đếm cục bộ cung cấp các ngưỡng; bộ nối mô tả ràng buộc.
  \item Thanh ghi tiến và lùi tách một cửa sổ cắt biên thành hai phần.
  \item Chứng minh tính đúng tách cửa sổ nội khối khỏi cửa sổ cắt biên.
  \item Độ rộng khối cân bằng chi phí thanh ghi với chi phí bộ nối.
\end{enumerate}
\end{summarybox}

\subsection*{Câu hỏi nghiên cứu}
\begin{enumerate}
  \item Dựng các bộ nối cho một cửa sổ AMK với \(k=2\) cắt biên thành hai phần.
  \item Tính số thanh ghi khi \(n=20,w=5,k=3\) và độ rộng khối bằng \(5\).
  \item Đề xuất hàm mục tiêu cho phân hoạch thích nghi dựa trên số cửa sổ cắt biên.
\end{enumerate}
